\section{Introduction}

\subsection{Decomposition theorem and inverse reduction}

To a simple Lie algebra $\g$ and a nilpotent element $f$, one can associate two interesting algebraic structures: 
\begin{enumerate}
	\item The \emph{finite W-algebra} $\U(\g, f)$: it is an associative algebra, constructed as the quantum Hamiltonian reduction of the universal enveloping algebra $\U(\g)$ \cite{kostant1978whittaker,premet2002special,gan2002quantization}.
	
	\item The \emph{affine W-algebra} $\W^k(\g, f)$: it is a vertex algebra, constructed as the chiral Hamiltonian reduction of the universal affine vertex algebra $\V^k(\g)$ \cite{feigin1990quantization,kac2003quantum,kac2004quantum}.
\end{enumerate}
The finite W-algebra is related to the affine W-algebra: the former is the associated \emph{Zhu algebra} of the latter \cite{desole2006finite}. Therefore, both their representation theories are related \cite{zhu1996modular}.

Losev introduced and studied an $\hbar$-adic analogue $\U(\g, f)_\hbar$ of the finite W-algebra \cite{losev2010quantized}, which is a complete $\C[\![\hbar]\!]$-algebra. He also established his celebrated \emph{decomposition theorem}, that is a $\hbar$-adic algebra isomorphism
\begin{equation} \label{equation:decomposition-theorem}
	\U(\g)_\hbar \cong A_\hbar \widehat\otimes_{\C[\![\hbar]\!]} \U(\g, f)_\hbar, 
\end{equation}
where $\U(\g)_\hbar$ is an $\hbar$-adic analogue of the enveloping algebra and $A_\hbar$ is a $\hbar$-adic Weyl algebra. Such a result is a crucial tool to connect the representation theories of $\g$ and the finite W-algebra $\U(\g, f)$. His proof is based on the geometry of the adjoint orbit of $f$ and its transverse slice $S_f$, called the \emph{Slodowy slice} of $f$.

For affine W-algebras, a similar kind of phenomenon has been intensively studied under the name of the \emph{inverse Hamiltonian reduction}. By definition, inverse reduction holds if there exists a vertex algebra embedding of the form
\begin{equation} \label{equation:inverse-reduction}
	\V^k(\g) \hooklongrightarrow  \Pi_f \otimes \W^k(\g, f), 
\end{equation}
where $\Pi_f$ is a certain localisation of a vertex algebraic analogue of a Weyl algebra. Such an embedding has been used to study weight modules of the simple quotient of $\V^k(\g)$ by using the representation theory of $\Pi_f$ and $\W^k(\g, f)$ \cite{adamovic2019realizations,adamovic2021realisation,adamovic2024relaxed,fasquel2025modularity}. 

Such inverse reductions have been established following different approaches:
\begin{itemize}
	\item When $\g$ is of small rank, inverse reduction can be constructed by using explicit generator and relations descriptions of W-algebras \cite{semikhatov1994inverting,adamovic2019realizations,adamovic2021realisation,adamovic2024relaxed}.
	
	\item By using the screening operator description of W-algebras \cite{genra2020screening}, inverse reduction has been established for any nilpotent $f$ when $\g$ is of type A \cite{fehily2023subregular,fehily2024inverse,fasquel2026connecting,fasquel2026quantum}.
\end{itemize}


To prove the decomposition theorem \eqref{equation:decomposition-theorem}, Losev uses the fact that the finite W-algebra $\U(\g, f)$ (and its $\hbar$-adic analogue) is a quantisation of the Poisson algebra $\C[S_f]$ of polynomial functions on the Slodowy slice $S_f$ \cite{premet2002special,gan2002quantization}. Such a relation also holds for the affine W-algebra: $\W^k(\g, f)$ carries a filtration \cite{li2004vertex,li2005abelianizing} such that the associated graded vertex algebra $\gr \W^k(\g, f)$ is commutative and isomorphic to $\C[\J_\infty S_f]$, the algebra of polynomial functions on the \emph{arc space} $\J_\infty S_f$ of the Slodowy slice \cite{desole2006finite,arakawa2015associated}. However, the approaches described above to establish inverse reduction \eqref{equation:inverse-reduction} do not make use of the Slodowy slice.

The Poisson structure of the Slodowy slice can be described by Hamiltonian reduction \cite{kostant1978whittaker,premet2002special,gan2002quantization}. There is a unipotent group $G_{\geqslant 1}$ that acts on the dual space $\g^\vee$, and a moment map $\mu \colon \g^\vee \rightarrow (\g_{\geqslant 1})^\vee$, such that $S_f$ is isomorphic to the Hamiltonian reduction $\mu^{-1}(0) / G_{\geqslant 1}$. In fact, there is a $G_{\geqslant 1}$-equivariant isomorphism
\begin{equation} \label{equation:gan-ginzburg} 
	G_{\geqslant 1} \times S_f \cong \mu^{-1}(0) 
\end{equation}
of varieties (see \Cref{theorem:gan-ginzburg-isomorphism}). 

This isomorphism is similar to the decomposition theorem \eqref{equation:decomposition-theorem} and the inverse reduction \eqref{equation:inverse-reduction}, with $S_f$ playing the role of the corresponding W-algebras. However, it is not the right geometric analogue because the dimension of the two varieties is too small and these varieties are not Poisson. The right analogue of \eqref{equation:decomposition-theorem} and \eqref{equation:inverse-reduction} must involve Poisson varieties of the same dimension as the Lie algebra $\g$. 

\subsection{A geometric approach: the affine Grassmannian}


Recently, a geometrical understanding of inverse reduction has been developed by using the affine Grassmannian \cite{kamnitzer2022hamiltonian,butson2025inverse,butson2025inverse-affine}. By the Mirkovi\'c--Vybornov correspondence \cite{mirkovic2022comparison}, Slodowy slices in type A are in correspondence with slices in the affine Grassmannian. A quantum analogue of this statement relates finite W-algebras in type A and truncated shifted Yangians \cite{webster2020quantum}. Such a correspondence is not established yet for affine W-algebras.

In \cite{kamnitzer2022hamiltonian,butson2025inverse}, the multiplication map of the affine Grassmannian is used to provide an inverse reduction statement for Slodowy slices and finite W-algebras, through the Mirkovi\'c--Vybornov correspondence. In particular, the multiplication map gives an analogue of the isomorphism \eqref{equation:gan-ginzburg} with the correct dimension: one side is an open subset of $\mathrm{T}^\vee G_{\geqslant 1} \times S_f$ (the product of the cotangent bundle of the unipotent group with the Slodowy slice), and the other side is an open subset of $\g^\vee$.

In \cite{butson2025inverse-affine}, Butson and Nair use the chiral quantisation theory they have developed in \cite{butson2026deformation} and the inverse reduction for Slodowy slices to establish inverse reduction \eqref{equation:inverse-reduction} in type A \cite[Section 6.3]{butson2025inverse-affine}. Note that only the existence of such an isomorphism is provided by the machinery of \cite{butson2026deformation}, but there is no explicit description of the map. The rough idea is to prove that both sides of the isomorphism are chiral quantisation of the same Poisson variety, and then use the classification of such quantisations to conclude that they must be isomorphic as vertex algebras. This argument is the vertex algebraic counterpart to Losev's approach for finite W-algebras \cite{losev2010quantized}.

\subsection{Our contribution}

In the present paper, we propose a new approach to inverse reduction that starts from the construction of the affine W-algebra as a chiral Hamiltonian reduction \cite{feigin1990quantization,kac2003quantum}. This reduction is concretely achieved by using the \emph{BRST (Becchi–Rouet–Stora–Tyutin) cohomology} functor (\Cref{section:brst}) $\Hgy^0_f$,
\[ \W^k(\g, f) \coloneqq \Hgy^0_f(\V^k(\g)). \]
Our construction is achieved in type A, for the minimal nilpotent orbit. For the rest of the paper, let $\g \coloneqq \sl_d$ be a simple Lie algebra of type A, $d \geqslant 2$ being an integer. Let $(e, h, f)$ be a minimal $\sl_2$-triple in $\g$.

The starting point of our construction is to establish an analogue of \eqref{equation:gan-ginzburg} between varieties of the correct dimension, without using the geometry of the affine Grassmannian. The $\sl_2$-triple equips $\g$ with a grading and $G_{\geqslant 1}$ is the unipotent group associated to the positive degree subalgebra $\g_{\geqslant 1}$. In \Cref{subsection:slodowy-quotient-open}, we construct a bigger group $K \ltimes G_{\geqslant 1}$, and its associated Lie algebra $\kfk \ltimes \g_{\geqslant 1}$ (see \Cref{lemma:afk}), such that the action of this group provides the desired isomorphism.

\begin{main}[\Cref{proposition:dilation-gan-ginzburg} and \Cref{theorem:decomposition-slodowy}] \label{main:geometric-setting}
	The group $K \ltimes G_{\geqslant 1}$ acts on the distinguished open set 
	\[ \g^\vee[e^{-1}] \coloneqq \Spec \big((\Sym \g)[e^{-1}]\big), \] 
	such that there is a $K \ltimes G_{\geqslant 1}$-equivariant isomorphism
	\[ (K \ltimes G_{\geqslant 1}) \times S_f \cong \g^\vee[e^{-1}] \]
	induced by the inclusion of the minimal Slodowy slice $S_f$ in $\g^\vee[e^{-1}]$.
\end{main}

By use of a ``quantisation commutes with reduction'' argument (\Cref{proposition:reduction-quantisation}), we can use \Cref{main:geometric-setting} to prove the main result of this paper.

\begin{main}[\Cref{theorem:affine-w-algebra-coset,theorem:decomposition-chiral}] \label{main:chiral-setting}
	One can construct a \emph{localisation} $\V^k(\g)[e^{-1}]$ of the universal affine vertex algebra $\V^k(\g)$ such that:
	\begin{enumerate}
		\item The affine W-algebra $\W^k(\g, f)$ is naturally isomorphic to a commutant subalgebra of $\V^k(\g)[e^{-1}]$:
		\[ \W^k(\g, f) \cong \{a \in \V^k(\g)[e^{-1}] \quad | \quad \text{$a$ commutes with $\V^k(\kfk \ltimes \g_{\geqslant 1})$} \}, \]
		where $\V^k(\kfk \ltimes \g_{\geqslant 1})$ is the vertex subalgebra of $\V^k(\g)$ generated by $\kfk \ltimes \g_{\geqslant 1}$.
		
		\item There is a vertex algebra isomorphism
		\begin{equation} \label{equation:main-decomposition-theorem}
			\V^k(\kfk \ltimes \g_{\geqslant 1})[e^{-1}] \otimes \W^k(\g, f) \overset{\sim}{\longrightarrow}  \V^k(\g)[e^{-1}] 
		\end{equation}
		given by the inclusion of the two mutually commuting vertex subalgebras $\V^k(\kfk \ltimes \g_{\geqslant 1})[e^{-1}]$ and $\W^k(\g, f)$ in $\V^k(\g)[e^{-1}]$.
	\end{enumerate}
\end{main}

The isomorphism \eqref{equation:main-decomposition-theorem} is an analogue of Losev's decomposition theorem, for non-$\hbar$-adic affine W-algebras. The theorem holds for any level $k$, no genericness is assumed. Note that, unlike the construction of \cite{butson2025inverse-affine}, the isomorphism that we construct is given by an explicit map. 

By \Cref{theorem:commutant-localisation-weyl}, the vertex subalgebra $	\V^k(\kfk \ltimes \g_{\geqslant 1})[e^{-1}]$ is in fact the localised Weyl vertex algebra $\Pi_f$ that appears in the literature on inverse reduction. \emph{Therefore, inverse reduction follows from our theorem in a straightforward way.} Indeed, localisation gives a vertex algebra embedding $\V^k(\g) \hookrightarrow \V^k(\g)[e^{-1}]$ (\Cref{subsection:localisation-vertex}). By composing this embedding with the inverse of the decomposition isomorphism \eqref{equation:main-decomposition-theorem}, we get the inverse reduction \eqref{equation:inverse-reduction}. 

The vertex algebra homomorphisms appearing in \Cref{main:chiral-setting} are not difficult to construct. The crucial point is to prove that they are isomorphisms. To do so, we use the increasing Li filtrations on vertex algebras (\Cref{subsection:vertex-algebras}). Taking the associated graded objects provide commutative algebras that correspond to polynomial functions on schemes:
\[ \gr \V^k(\g) \cong \C[\J_\infty \g^\vee],  \quad \gr \big(\V^k(\g)[e^{-1}]\big) \cong \C[\J_\infty \g^\vee][e^{-1}], \quad   \gr \W^k(\g, f) \cong \C[\J_\infty S_f], \]
where $\J_\infty X$ denotes the arc scheme of a scheme $X$ (\Cref{subsection:arc-schemes}).


Finally, \Cref{main:geometric-setting} also leads to a decomposition theorem and inverse reduction statement for the minimal finite W-algebra.

\begin{main}[\Cref{theorem:finite-w-algebra-coset,theorem:decomposition-quantum}] \label{main:quantum-setting}
	The minimal finite W-algebra $\U(\g, f)$ is naturally isomorphic to the commutant subalgebra
	\[ (\U(\g)[e^{-1}])^{\kfk \ltimes \g_{\geqslant 1}} \coloneqq \{a \in \U(\g)[e^{-1}] \quad | \quad \text{$a$ commutes with $\kfk \ltimes \g_{\geqslant 1}$} \}. \]
	Moreover, there is an algebra isomorphism
	\begin{equation} \label{equation:main-decomposition-theorem-finite}
		\U(\kfk \ltimes \g_{\geqslant 1})[e^{-1}] \otimes \U(\g, f) \overset{\sim}{\longrightarrow}  \U(\g)[e^{-1}]. 
	\end{equation}
\end{main}

\subsection{The $\sl_2$ example}

We illustrate our result in the easiest case, when $\g = \sl_2$ and
\[ e = \left(\begin{smallmatrix}
	0 & 1 \\
	0 & 0
\end{smallmatrix}\right), \quad h = \left(\begin{smallmatrix}
	1 & 0 \\
	0 & -1
\end{smallmatrix}\right) \quad \text{and} \quad f = \left(\begin{smallmatrix}
	0 & 0 \\
	1 & 0
\end{smallmatrix}\right). \]

\subsubsection{The geometric setting}

The symmetric algebra $\Sym \sl_2$ has a Poisson structure given by
\[ \{h, e\} = 2 e, \quad \{e, f\} = h \quad \text{and} \quad \{h, f\} = - 2 f. \]

In this case the Slodowy slice is given by the explicit description:
\[ S_f \coloneqq \Big\{\begin{pmatrix}
	0 & \star \\
	1 & 0
\end{pmatrix}\Big\}. \]
By the work of Kostant \cite{kostant1978whittaker}, the Poisson centre $\mathrm{Z}(\Sym \sl_2)$ is generated by the Casimir element
\[ c \coloneqq \frac{1}{2} h^2 + 2 e f \] 
and isomorphic to $\C[S_f]$.

There is an obvious Poisson algebra homomorphism
\[ \Sym(\C h \oplus \C e) \otimes \mathrm{Z}(\Sym \sl_2) \longrightarrow \Sym \sl_2. \]
Note that if we allow $e$ to be invertible, then we can write
\[ f = e^{-1} \frac{1}{2}(c - \frac{1}{2} h^2). \]
Therefore we have a Poisson algebra isomorphism
\[ \Sym(\C h \oplus \C e)[e^{-1}] \otimes \mathrm{Z}(\Sym \sl_2) \longrightarrow (\Sym \sl_2)[e^{-1}]. \]

\subsubsection{The quantum setting}

The enveloping algebra $\U(\sl_2)$ has commutations relations given by
\[ [h, e] = 2 e, \quad [e, f] = h \quad \text{and} \quad [h, f] = - 2 f. \]

By the work of Kostant \cite{kostant1978whittaker}, the finite W-algebra $\U(\sl_2, f)$ is isomorphic to the centre $\mathrm{Z}(\U(\sl_2))$, that is generated by the Casimir element
\[ \widetilde c \coloneqq \frac{1}{2} h^2 + 2 e f - h. \] 

There is an obvious algebra homomorphism
\[ \U(\C h \oplus \C e) \otimes \mathrm{Z}(\U(\sl_2)) \longrightarrow \U(\sl_2). \]
Note that if we allow $e$ to be invertible, then we can write
\[ f = e^{-1} \frac{1}{2}(\widetilde{c} - \frac{1}{2} h^2 + h). \]
Therefore we have an algebra isomorphism
\[ \U(\C h \oplus \C e)[e^{-1}] \otimes \mathrm{Z}(\U(\sl_2)) \cong \U(\sl_2)[e^{-1}]. \]

\subsubsection{The chiral setting at critical level}

The universal affine vertex algebra $\V^k(\sl_2)$ has operator product expansions given by
\[ \begin{multlined}
	h(z_1) h(z_2) \sim \frac{1}{(z_1 - z_2)^2} 2 k \id, \quad h(z_1) e(z_2) \sim \frac{1}{z_1 - z_2} 2 e(z_2) , \\ e(z_1) f(z_2) \sim \frac{1}{z_1 - z_2} h(z_2) + \frac{1}{(z_1 - z_2)^2} k \id \\ \text{and} \quad h(z_1) f(z_2) \sim - \frac{1}{z_1 - z_2} 2 f(z_2).
\end{multlined} \]
To simplify the exposition, we consider only the critical level $k = -2$. For the general case, see \Cref{section:sl2-example}.

By the work of Feigin--Frenkel \cite{feigin1992affine}, at the critical level, the affine W-algebra $\W^{-2}(\sl_2, f)$ is isomorphic to the centre of $\V^{-2}(\sl_2)$ and is generated by
\[ \widehat c \coloneqq \frac{1}{2} \NO{h h} + \NO{e f} + \NO{f e} = \frac{1}{2} \NO{h h} + 2 \NO{e f} - \partial h, \]
where $\partial$ denotes the translation operator and $\NO{}$ denotes the normally ordered product. 

There is an obvious vertex algebra homomorphism
\[ \V^{-2}(\C h \oplus \C e) \otimes \W^{-2}(\sl_2, f) \longrightarrow \V^{-2}(\sl_2). \]
Note that if we allow $e$ to be invertible, then we can write
\[ \NO{e^{-1} (\NO{e f})} = \NO{e^{-1} \frac{1}{2}(\widehat{c} - \frac{1}{2} \NO{h h} + \partial h)}. \]
The non-associativity of the normally ordered product is controlled by the following Borcherds identity:
\[ \begin{multlined}
	f = {(\NO{e^{-1} e}) f} = \NO{e^{-1} (\NO{e f})} + P, 
	\\ \text{where} \quad P \coloneqq \sum_{i \geqslant 0} \big(e^{-1}{}_{(-2-i)}(e_{(i)} f) + e{}_{(-2-i)}(e^{-1}{}_{(i)} f)\big), 
\end{multlined}  \]
with the notation $\bullet_{(i)}\bullet$ for the $i$-products in the vertex algebra.

The operator product expansion of $e$ and $f$ says that $e_{(i)}f$ belongs to the subalgebra $\V^{-2}(\C h \oplus \C e)$ for any nonnegative integer $i$. Then, \cite[Formula (3.1)]{arakawa2026faithful} implies that $e^{-1}{}_{(i)}f$ belongs to $\V^{-2}(\C h \oplus \C e)[e^{-1}]$ for any nonnegative integer $i$. Therefore, $P$ belongs to $\V^{-2}(\C h \oplus \C e)[e^{-1}]$ and
\[ f = \NO{e^{-1} \frac{1}{2}(\widehat{c} - \frac{1}{2} \NO{h h} + \partial h)} + P. \]
See \Cref{proposition:computation-f} for an explicit formula.

Therefore we get a vertex algebra isomorphism
\[ \V^{-2}(\C h \oplus \C e)[e^{-1}] \otimes \W^{-2}(\sl_2, f) \cong \V^{-2}(\sl_2)[e^{-1}]. \]


\subsection{Future work and applications}

\subsubsection{Reduction by stages in type A}

In a work in progress \cite{juillard2027}, we extend our construction to inverse reduction by stages in type A. By using our previous work with Genra \cite{genra2024reduction,genra2025reduction}, for two adjacent hook-type nilpotent elements $f_1$ and $f_2$ in $\g = \sl_d$ such that the orbit of $f_2$ is bigger than the orbit of $f_1$, we construct an explicit vertex algebra isomorphism
\[ \Pi_{f_1, f_2} \otimes \W^k(\g, f_2) \cong \W^k(\g, f_1)[e_0{}^{-1}], \]
where $e_0$ is a localisable element in $\W^k(\g, f_1)$ and $\Pi_{f_1, f_2}$ is a localised Weyl vertex algebra. By induction, we can deduce a decomposition theorem for any hook-type nilpotent element $f$ in $\g$:
\[ \Pi_{f} \otimes \W^k(\g, f) \cong \V^k(\g)[S], \]
where $S$ is a localisable subset in $\V^k(\g)$ and $\Pi_{f}$ is a localised Weyl vertex algebra.

\subsubsection{Minimal W-algebra in any type}

If $\g$ is not of type A, \Cref{main:geometric-setting} no longer holds: we need to replace the open set $\g^\vee[e^{-1}]$ by a suitable étale cover. On the noncommutative side, this means that we need to replace the localisation embedding $\V^k(\g) \hookrightarrow \V^k(\g)[e^{-1}]$ by a more complicated extension of $\V^k(\g)$ to get the right analogue of \Cref{main:chiral-setting} for a general $\g$. 

With Christopher Raymond, we aim to establish a decomposition theorem for any simple Lie algebra $\g$ of rank $2$, that is to say of type $\mathrm{A}_2$, $\mathrm{B}_2$ or $\mathrm{G}_2$. Note that an inverse reduction in type $\mathrm{A}_2$ is already known and is the one retrieved by \Cref{main:chiral-setting}. However, replacing $\V^k(\sl_3) \hookrightarrow \V^k(\sl_3)[e^{-1}]$ by a more complicated extension of $\V^k(\sl_3)$ should lead to another inverse reduction. 

As mentioned earlier, inverse reduction has proved useful to study weight modules over the simple quotient of $\V^k(\g)$ \cite{adamovic2019realizations,adamovic2021realisation,adamovic2024relaxed,fasquel2025modularity}. We want to pursue these studies in these new cases. 

\subsection{Organisation of the paper}

In \Cref{section:geometric-setting}, we recall the construction of the Slodowy slice as Hamiltonian reduction and establish the decomposition theorem for the minimal Slodowy slice (\Cref{theorem:decomposition-slodowy}). An important result for the rest of the paper is \Cref{proposition:dilation-gan-ginzburg}, that is an analogue of the Gan--Ginzburg isomorphism (\Cref{theorem:gan-ginzburg-isomorphism}) for an open set in $\g^\vee$. 

In \Cref{section:localisation-filtration}, we collect technical results on filtrations on noncommutative algebras and localisation in the noncommutative setting. We also give some results to compute invariant subspaces of filtered modules for an algebraic group.

The results of the two previous sections are applied in \Cref{section:quantum-setting} to prove that the minimal finite W-algebra is an invariant subalgebra of the localised enveloping algebra (\Cref{theorem:finite-w-algebra-coset}). Then, we establish the decomposition theorem for the minimal finite W-algebra (\Cref{theorem:decomposition-quantum}).

In \Cref{section:vertex-algebra}, we collect results about infinite type objects: vertex algebras and arc spaces. We also recall the Borisov localisation theory of vertex algebras and its compatibility with the Li filtration. We give some technical results on computing invariant subspaces of filtered modules for the projective limit of a system of algebraic groups. 

In \Cref{section:brst}, we recall the BRST construction of chiral Hamiltonian reduction and we give a sufficient condition to prove that ``quantisation commutes with reduction'' (\Cref{proposition:reduction-quantisation}).

In \Cref{section:chiral-setting}, we prove that the minimal affine W-algebra is an invariant subalgebra of the localised universal vertex algebra (\Cref{theorem:affine-w-algebra-coset}). Then, we establish the decomposition theorem for the minimal affine W-algebra (\Cref{theorem:decomposition-chiral}). Finally, we prove that the subalgebra of the localisation commuting with the W-algebra is a localised Weyl vertex algebra (\Cref{theorem:commutant-localisation-weyl}).

In \Cref{section:sl2-example}, we give the example of our construction, in the vertex-algebraic setting, for $\g = \sl_2$. In particular, we make the computations at any level.

\subsection{Notations and conventions}

Unless otherwise stated, all objects (vector spaces, algebras, schemes, tensor products, ...) are defined over the field of complex numbers $\C$. Denote by $\Z$ the ring of integers. Denote by $\delta_{a=b}$ the Kronecker symbol that is equal to $1$ if $a=b$ holds, $0$ otherwise. 

If $V$ is a vector space, denote by $\Sym V$ its symmetric algebra, $\operatorname{Alt} V$ its exterior algebra and $V^\vee$ its dual vector space.

For any affine scheme $X = \Spec R$, denote by $\C[X] \coloneqq R$ its coordinate ring. If $\phi : X \rightarrow Y$ is a homomorphism of affine schemes, the comorphism is denoted by~$\phi^* : \C[Y] \rightarrow \C[X]$. If $R$ is a $\C$-algebra and $X$ is a scheme, then $X(R)$ denotes the set of $R$-points of $X$, that is to say the set of unital $\C$-algebra homomorphisms from $\C[X]$ to $R$.

If $S$ is a multiplicative subset in $R$, denote by $R[S^{-1}]$ the corresponding localisation. Denote by $X[S^{-1}] \coloneqq \Spec R[S^{-1}]$ the corresponding open affine subscheme. If $S$ is generated by one element $F$ in $R$, then we adopt the notation $R[F^{-1}] \coloneqq R[S^{-1}]$ and $X[F^{-1}] \coloneqq \Spec R[F^{-1}]$.  

If $X$ is a variety (reduced affine scheme of finite type), the notation $x \in X$ means that $x$ is a $\C$-point of $X$.  A linear algebraic group is an affine variety with a group scheme structure. The multiplicative group is denoted by~$\Gm \coloneqq \Spec \C[t, t^{-1}]$ and is identified with its $\C$-point group of invertible complex numbers $\C^\times$. The additive group is denoted by~$\Ga \coloneqq \Spec \C[t]$ and is identified with its $\C$-point group of complex numbers $\C$.

If $G$ is a linear algebraic group, its Lie algebra is denoted by $\g$. Actions of an algebraic group on a scheme or a vector space are assumed to be algebraic, that is to say given by a co-action \cite{milne2017algebraic}. 

If both $G$ and $\g$ act on a vector space $V$, for any vector $v$ in $V$ set
$$ G^v \coloneqq \{g \in G \quad | \quad g \cdot v = v \} \quad \text{and} \quad \g^v \coloneqq \{x \in \g \quad | \quad x \cdot v = 0 \}. $$
Note that $\g^v$ is the Lie algebra of $G^v$. The subspaces of invariants are denoted by
\[ \begin{aligned}
	V^G & \coloneqq \{v \in V \quad | \quad \text{for all $g \in G$},~ g \cdot v = v \} \\
	\text{and} \quad V^\g & \coloneqq \{v \in V \quad | \quad \text{for all $x \in \g$},~ x \cdot v = 0 \}
\end{aligned}\]

If $G$ acts on an affine scheme $X \coloneqq \Spec R$, the affine categorical quotient is the affine scheme
\[ X \git G \coloneqq \Spec R^G. \]


\subsection{Acknowledgements}

The idea of this project arose in a discussion with Naoki Genra on the $\g = \sl_2$ case. I am very grateful to him for many years of stimulating discussions and collaboration. For the same reason, I am grateful to Christopher Raymond. In particular, our conversations were essential to achieve the proof of \Cref{theorem:commutant-localisation-weyl} and the computations of \Cref{section:sl2-example}.

I thank Anne Moreau and Sven M\"oller for their advice during the preparation of this paper and for their careful proofreading. I thank Tomoyuki Arakawa, Gwyn Bellamy, Thomas Creutzig and David Ridout for enlightening discussions. I also thank Tomoyuki Arakawa for communicating the results of \cite{arakawa2026faithful} to me. I thank Felipe Ruiz and Tanguy Vernet for their feedback on the introduction.

The $\g = \sl_2$ case has been computed independently by Tomoyuki Arakawa and Thomas Creutzig. 

This research was supported through the Emmy Noether Programme by
the \emph{Deutsche Forschungsgemeinschaft} (project number 460925688).

\subsection{Publication overlap}

During the proofreading of the present paper, we noticed the prepublication of Genqiang Liu \cite{liu2026root}. The main result of Liu is the same as our \Cref{main:quantum-setting}, but for $\g$ being of any type. 

\subsection{AI disclosure}

No generative or agentive AI model was used to carry out this research. A first version of this paper was written without AI assistance. \emph{DeepSeek~V4.1~Flash} was used by the author to produce a report on grammar, reference and notation mistakes. All the changes were then implemented manually by the author to produce the current version of the paper.